\chapter{Bilangan Cacah}\label{ch:g6:wholes}

Semua hal dalam matematika bermula dari membilang. Bab ini meninjau
kembali cara bilangan cacah ditulis, dibandingkan dan digabungkan ---
serta memeriksa dengan teliti pembagian, operasi yang memberikan \omterm{def:g3:division:remainder}{hasil
bagi} \emph{dan} \omterm{def:g3:division:remainder}{sisa} sekaligus.

\section{Menulis dan membandingkan bilangan cacah}

\begin{definition}[Nilai tempat]\label{def:g6:wholes:place}
Cara kita menulis bilangan memakai sepuluh angka ($0$ sampai $9$), dan
\emph{nilai sebuah angka bergantung pada tempatnya}\index{nilai tempat}:
pada $5\,208$, angka $5$ membilang ribuan, angka $2$ membilang ratusan,
angka $0$ membilang puluhan (tidak ada) dan angka $8$ membilang satuan:
\[
5\,208 = 5 \times 1000 + 2 \times 100 + 0 \times 10 + 8 .
\]
\end{definition}

\begin{omfigure}
\begin{tikzpicture}[scale=0.9]
  \draw (0,0) grid[xstep=2.3, ystep=0.8] (9.2,1.6);
  \node[font=\small] at (1.15,1.2) {ribuan};
  \node[font=\small] at (3.45,1.2) {ratusan};
  \node[font=\small] at (5.75,1.2) {puluhan};
  \node[font=\small] at (8.05,1.2) {satuan};
  \node[font=\large, omDef] at (1.15,0.4) {$5$};
  \node[font=\large, omDef] at (3.45,0.4) {$2$};
  \node[font=\large, omDef] at (5.75,0.4) {$0$};
  \node[font=\large, omDef] at (8.05,0.4) {$8$};
\end{tikzpicture}

{\small Tabel \omterm{def:g6:wholes:place}{nilai tempat} dari $5\,208$. Angka $0$ itu penting: tanpa
dia, angka $5$ dan $2$ akan tergelincir ke kolom yang salah dan
bilangannya terbaca $528$.}
\end{omfigure}

\begin{method}[Membandingkan dua bilangan cacah]\label{met:g6:wholes:compare}
\begin{enumerate}
  \item Bilangan yang angkanya lebih banyak adalah yang lebih besar
        ($1\,203 > 989$).
  \item Jika banyak angkanya sama, bandingkan angka demi angka dari
        \emph{kiri}; perbedaan pertama yang menentukan:
        $6\,742 > 6\,715$ karena pada tempat puluhan, $4 > 1$.
\end{enumerate}
Lambangnya adalah $<$ (``lebih kecil daripada'') dan $>$ (``lebih
besar daripada''); ujung yang kecil menunjuk ke bilangan yang kecil.
\end{method}

\begin{example}\label{ex:g6:wholes:order}
Urutkan dari yang terkecil ke yang terbesar: $908$; $89$; $1\,001$;
$980$. Mulailah dari banyaknya angka: $89$ (dua angka) di depan,
$1\,001$ (empat angka) di belakang. Antara $908$ dan $980$: angka
ratusannya sama, $9$, lalu puluhannya $0 < 8$, jadi $908 < 980$.
Jawaban akhirnya:
\[
89 < 908 < 980 < 1\,001 .
\]
\end{example}

\begin{omfigure}
\begin{tikzpicture}[scale=1.05]
  \draw[-Stealth] (-0.4,0) -- (10.8,0);
  \foreach \x in {0,1,...,10} \draw (\x,-0.08) -- (\x,0.08);
  \foreach \x/\l in {0/0, 2/200, 4/400, 6/600, 8/800, 10/1000}
    \node[below, font=\small] at (\x,-0.1) {$\l$};
  \fill[omThm] (0.89,0) circle (2pt) node[above, font=\small] {$89$};
  \fill[omDef] (9.08,0) circle (2pt) node[above, font=\small] {$908$};
  \fill[omMeth] (9.8,0) circle (2pt) node[above right=-2pt, font=\small] {$980$};
\end{tikzpicture}

{\small Bilangan tinggal pada sebuah garis: makin kecil berarti makin
ke kiri. Setiap tanda di sini bernilai $100$.}
\end{omfigure}

\section{Penjumlahan, pengurangan, perkalian}

\begin{example}[Perhitungan bersusun]\label{ex:g6:wholes:columns}
Untuk menjumlahkan atau mengurangkan, luruskan satuan di bawah satuan,
puluhan di bawah puluhan, lalu kerjakan dari kanan ke kiri sambil
menyimpan bila perlu:
\[
\begin{array}{r}
4\,5\,7 \\
+\ 2\,8\,6 \\
\hline
7\,4\,3
\end{array}
\qquad\qquad
\begin{array}{r}
6\,0\,3 \\
-\ 1\,4\,7 \\
\hline
4\,5\,6
\end{array}
\]
Pada penjumlahannya: $7 + 6 = 13$, tulis $3$, simpan $1$; lalu
$5 + 8 + 1 = 14$, tulis $4$, simpan $1$; lalu $4 + 2 + 1 = 7$.
Periksa pengurangannya dengan menjumlahkan kembali: $456 + 147 = 603$.
\end{example}

\begin{remark}[Kosakata]
Hasil sebuah penjumlahan disebut \emph{jumlah}\index{jumlah}; hasil
pengurangan, \emph{selisih}\index{selisih}; hasil perkalian,
\emph{hasil kali}\index{hasil kali}; hasil pembagian,
\emph{hasil bagi}\index{hasil bagi}. Bilangan yang digabungkan disebut
\emph{suku} (untuk $+$, $-$) atau \emph{faktor} (untuk $\times$).
\end{remark}

\begin{example}[Menghitung dengan cerdik]\label{ex:g6:wholes:smart}
Menukar urutan suku atau faktor sering menghemat pekerjaan:
\[
17 + 58 + 3 = (17 + 3) + 58 = 20 + 58 = 78 ,
\]
\[
25 \times 17 \times 4 = (25 \times 4) \times 17 = 100 \times 17 = 1700 .
\]
\end{example}

\section{Pembagian}

\begin{theorem}[Pembagian Euklides]\label{thm:g6:wholes:division}
Diberikan dua bilangan cacah $a$ (\emph{bilangan yang dibagi}) dan
$b \neq 0$ (\emph{\omterm{def:g5:division:multiple}{pembagi}}), ada tepat satu cara untuk menulis
\[
a = b \times q + r
\qquad\text{dengan}\qquad
0 \leq r < b .
\]
Bilangan $q$ disebut \emph{hasil bagi}\index{hasil bagi} dan $r$
disebut \emph{sisa}\index{sisa} pembagian $a$ oleh $b$.
\end{theorem}
\admitted

\begin{example}\label{ex:g6:wholes:division}
Bagilah $137$ dengan $4$, selangkah demi selangkah.
\begin{enumerate}
  \item Berapa kali $4$ masuk ke dalam $13$? Tiga kali
        ($4 \times 3 = 12$), \omterm{def:g3:division:remainder}{sisanya} $13 - 12 = 1$.
  \item Turunkan angka $7$: berapa kali $4$ masuk ke dalam $17$?
        Empat kali ($4 \times 4 = 16$), \omterm{def:g3:division:remainder}{sisanya} $1$.
\end{enumerate}
Jadi $q = 34$ dan $r = 1$:
\[
137 = 4 \times 34 + 1, \qquad\text{dan memang } 0 \leq 1 < 4 .
\]
Periksa: $4 \times 34 = 136$, dan $136 + 1 = 137$.
\end{example}

\begin{omfigure}
\begin{tikzpicture}[scale=0.95]
  \draw[-Stealth] (-0.3,0) -- (11.2,0);
  \foreach \i in {0,1,...,34} \draw ({\i*0.32},-0.06) -- ({\i*0.32},0.06);
  \foreach \i in {0,10,20,30}
    \node[below, font=\footnotesize] at ({\i*0.32},-0.08) {$4\times\i$};
  \draw[omDef, very thick] (0,0.14) -- (10.88,0.14);
  \draw[omThm, very thick] (10.88,0.14) -- (10.96,0.14);
  \fill (10.96,0) circle (1.8pt) node[above, font=\small] {$137$};
  \node[omDef, font=\small] at (5.4,0.55) {$34$ kelompok penuh berisi $4$};
\end{tikzpicture}

{\small Membagi $137$ dengan $4$: kamu dapat memuat $34$ kelompok
penuh berisi $4$ (mencapai $136$), dan $1$ tersisa --- \omterm{def:g3:division:remainder}{sisanya} selalu
lebih kecil daripada \omterm{def:g5:division:multiple}{pembaginya}.}
\end{omfigure}

\begin{definition}[Habis dibagi]\label{def:g6:wholes:divisible}
Ketika \omterm{def:g3:division:remainder}{sisanya} $0$, kita mengatakan bahwa $a$ \emph{habis
dibagi}\index{habis dibagi} oleh $b$: misalnya $135$ habis dibagi $5$
karena $135 = 5 \times 27 + 0$.
\end{definition}

\begin{method}[Soal cerita dengan pembagian]\label{met:g6:wholes:problems}
Setelah menghitung $a = bq + r$, kembalilah selalu ke pertanyaannya:
\begin{enumerate}
  \item ``Berapa kotak/regu/halaman yang penuh?'' --- jawabannya
        adalah \omterm{def:g3:division:remainder}{hasil bagi} $q$;
  \item ``Berapa yang tersisa?'' --- jawabannya adalah \omterm{def:g3:division:remainder}{sisa} $r$;
  \item ``Berapa kotak yang diperlukan untuk semuanya?'' --- kalau
        $r > 0$, satu \emph{lebih} daripada \omterm{def:g3:division:remainder}{hasil baginya}: $q + 1$.
\end{enumerate}
\end{method}

\begin{example}\label{ex:g6:wholes:bus}
$137$ murid pergi bertamasya dengan bus berisi $40$ tempat duduk.
Pembagiannya: $137 = 40 \times 3 + 17$. Tiga bus mengangkut $120$
murid, dan $17$ murid tersisa --- mereka juga perlu bus! Jadi $4$ bus
diperlukan, walaupun \omterm{def:g3:division:remainder}{hasil baginya} $3$.
\end{example}

\section{Latihan}

\begin{exercise}[$\star$]\label{exo:g6:wholes:1}
Tulislah dengan angka: ``tiga ribu empat puluh tujuh''; ``dua puluh
ribu lima ratus''; ``seratus empat''. Lalu tulislah $60\,013$ dengan kata.
\end{exercise}

\begin{exercise}[$\star$]\label{exo:g6:wholes:2}
Pada bilangan $47\,209$: berapa angka ratusannya? Angka satuannya?
Apa yang dibilang angka $4$? Tulislah bilangan itu sebagai \omterm{def:g1:addition:def}{jumlah}
\omterm{def:g5:division:multiple}{kelipatan} $1$, $10$, $100$, \dots\ seperti pada
\cref{def:g6:wholes:place}.
\end{exercise}

\begin{exercise}[$\star$]\label{exo:g6:wholes:3}
Salin dan lengkapi dengan $<$ atau $>$:
\[
507 \;?\; 570, \qquad
1\,099 \;?\; 989, \qquad
23\,456 \;?\; 23\,465, \qquad
9\,999 \;?\; 10\,000 .
\]
\end{exercise}

\begin{exercise}[$\star$]\label{exo:g6:wholes:4}
Urutkan dari yang terkecil ke yang terbesar: $2\,043$; $978$; $2\,304$; $2\,034$;
$1\,001$.
\end{exercise}

\begin{exercise}[$\star$]\label{exo:g6:wholes:5}
Hitunglah secara bersusun: $348 + 576$; \quad $805 - 268$; \quad
$307 \times 26$. Periksalah pengurangannya dengan penjumlahan.
\end{exercise}

\begin{exercise}[$\star$]\label{exo:g6:wholes:6}
Hitunglah dengan cerdik, dengan mengelompokkan suku atau faktornya:
\[
38 + 45 + 12 + 5, \qquad
2 \times 83 \times 50, \qquad
4 \times 78 \times 25 .
\]
\end{exercise}

\begin{exercise}[$\star$]\label{exo:g6:wholes:7}
Untuk setiap pembagian, sebutkan \omterm{def:g3:division:remainder}{hasil bagi} dan \omterm{def:g3:division:remainder}{sisanya}, lalu
tulislah kesamaan $a = b \times q + r$:
\[
95 \div 7, \qquad
230 \div 6, \qquad
504 \div 8 .
\]
\end{exercise}

\begin{exercise}[$\star$]\label{exo:g6:wholes:8}
Tanpa membagi, jelaskan mengapa kesamaan $173 = 8 \times 20 + 13$
\emph{bukan} pembagian Euklides $173$ oleh $8$, lalu carilah \omterm{def:g3:division:remainder}{hasil
bagi} dan \omterm{def:g3:division:remainder}{sisa} yang benar.
\end{exercise}

\begin{exercise}[$\star\star$]\label{exo:g6:wholes:9}
Telur dikemas dalam kotak berisi $12$.
\begin{enumerate}
  \item Sebuah peternakan punya $2\,000$ telur. Berapa kotak penuh
        yang dapat diisinya, dan berapa telur yang tersisa?
  \item Sebuah toko roti memerlukan $150$ telur. Berapa kotak yang
        harus dibelinya?
\end{enumerate}
\end{exercise}

\begin{exercise}[$\star\star$]\label{exo:g6:wholes:10}
Hari ini hari Selasa. Hari apa nanti dalam $100$ hari? (Satu minggu
ada $7$ hari: bagilah lalu pikirkan \omterm{def:g3:division:remainder}{sisanya}.)
\end{exercise}

\begin{exercise}[$\star\star$]\label{exo:g6:wholes:11}
Pada sebuah pembagian oleh $9$, \omterm{def:g3:division:remainder}{hasil baginya} $56$ dan \omterm{def:g3:division:remainder}{sisanya}
sebesar mungkin. Berapa bilangan yang dibagi itu?
\end{exercise}

\begin{exercise}[$\star\star\star$]\label{exo:g6:wholes:12}
Aku bilangan tiga angka. Angka ratusanku dua kali angka satuanku,
angka puluhanku $0$, dan \omterm{def:g1:addition:def}{jumlah} angkaku $9$. Siapakah aku?
(Bernalarlah selangkah demi selangkah, lalu periksa.)
\end{exercise}

\section{Soal: Gauss kecil dan jumlah seratus bilangan
yang pertama}

\begin{problem}[{Soal akhir pekan --- menjumlahkan $1 + 2 + 3 + \dots +
100$ dalam sepuluh detik, dan bilangan segitiga}]
\label{pb:g6:wholes:1}
Sebuah cerita yang terkenal: untuk menyibukkan kelasnya, seorang guru
meminta murid-muridnya menjumlahkan semua bilangan dari $1$ sampai
$100$. Seorang anak laki-laki menulis $5\,050$ di batu tulisnya
hampir seketika --- Carl Friedrich Gauss, umur sembilan tahun, yang
kelak tumbuh menjadi salah satu di antara matematikawan terbesar
sepanjang masa. Rahasianya tidak lain adalah \emph{menghitung dengan cerdik}
(\cref{ex:g6:wholes:smart}): mengelompokkan sukunya sehingga setiap
kelompok menjadi mudah. Soal ini menemukan kembali triknya, membuatnya
berhasil dalam segala keadaan, dan memperlihatkan betapa sering \omterm{def:g1:addition:def}{jumlah}
semacam itu muncul dalam kehidupan sehari-hari.

\medskip
\noindent\textbf{Bagian I --- Pasangan Gauss.}
\begin{enumerate}
  \item Hitunglah $1 + 2 + 3 + 4 + 5$ dengan mengelompokkan sukunya
        secara cerdik.
  \item Untuk $1 + 2 + 3 + \dots + 10$: pasangkan suku pertama dengan
        suku terakhir ($1 + 10$), suku kedua dengan suku kedua dari
        belakang ($2 + 9$), dan seterusnya. Berapa nilai setiap
        pasangan? Ada berapa pasangan? Selesaikan perhitungannya.
  \item Cara yang sama untuk $1 + 2 + \dots + 20$: berapa nilai
        setiap pasangan, ada berapa pasangan, berapa \omterm{def:g1:addition:def}{jumlahnya}?
  \item Sekarang \omterm{def:g1:addition:def}{jumlah} sang guru, $1 + 2 + \dots + 100$: jelaskan
        mengapa pasangan $1 + 100$, $2 + 99$, $3 + 98, \dots$
        semuanya bernilai sama, dan mengapa banyaknya tepat $50$,
        dengan pasangan terakhir $50 + 51$.
  \item Tulislah jawaban Gauss sebagai satu perkalian saja, lalu
        hitunglah.
\end{enumerate}

\medskip
\noindent\textbf{Bagian II --- Trik yang tidak pernah gagal.}
\begin{enumerate}[resume]
  \item Cobalah pemasangan itu pada $1 + 2 + \dots + 9$: kali ini
        satu bilangan tertinggal tanpa pasangan. Yang mana?
        Pasangkan di sekelilingnya lalu hitunglah \omterm{def:g1:addition:def}{jumlahnya}.
  \item Inilah bentuk trik yang berhasil baik ketika banyaknya suku
        \omterm{def:g3:numbers:evenodd}{genap} maupun \omterm{def:g3:numbers:evenodd}{ganjil}: tulislah \omterm{def:g1:addition:def}{jumlahnya} dua kali, sekali maju
        dan sekali mundur, satu baris di bawah yang lain,
        \[
        \begin{array}{ccccccc}
        1 & + & 2 & + & \dots & + & 9 \\
        9 & + & 8 & + & \dots & + & 1
        \end{array}
        \]
        lalu jumlahkan \emph{kolom demi kolom}. Berapa nilai setiap
        kolom? Ada berapa kolom? Jelaskan mengapa \omterm{def:g1:addition:def}{jumlah} kedua baris
        itu $9 \times 10 = 90$, lalu simpulkan bahwa
        $1 + 2 + \dots + 9 = 45$ --- jawaban yang sama seperti pada
        pertanyaan 6.
  \item Pakailah trik dua baris itu untuk menghitung
        $1 + 2 + \dots + 50$, lalu $1 + 2 + \dots + 200$.
  \item Hitunglah \omterm{def:g1:addition:def}{jumlah} seratus bilangan \emph{\omterm{def:g3:numbers:evenodd}{genap}} yang pertama,
        $2 + 4 + 6 + \dots + 200$. (Bandingkan setiap sukunya dengan
        suku pada \omterm{def:g1:addition:def}{jumlah} Gauss: tidak perlu pemasangan yang baru.)
  \item Hitunglah $101 + 102 + \dots + 200$, dengan memakai dua
        \omterm{def:g1:addition:def}{jumlah} yang sudah kamu ketahui dan satu pengurangan.
\end{enumerate}

\medskip
\noindent\textbf{Bagian III --- \omterm{pb:g6:wholes:1}{Bilangan segitiga} di mana-mana.}
\omterm{def:g1:addition:def}{Jumlah} $1$, $3$, $6$, $10$, $15, \dots$ dari deret
$1 + 2 + \dots + n$ disebut \emph{bilangan
segitiga}\index{bilangan segitiga}, karena bilangan itu membilang
benda yang ditumpuk membentuk segitiga.
\begin{enumerate}[resume]
  \item Delapan orang teman bertemu, dan masing-masing berjabat
        tangan tepat sekali dengan setiap temannya yang lain.
        Jelaskan mengapa banyaknya jabat tangan adalah
        $7 + 6 + 5 + 4 + 3 + 2 + 1$, lalu hitunglah.
  \item Seorang pedagang menyusun \omterm{def:g5:solids:def}{limas} segitiga dari kaleng: $12$
        kaleng di baris paling bawah, $11$ di baris berikutnya, dan
        seterusnya sampai satu kaleng di puncaknya. Berapa kaleng
        yang ia perlukan? (Pakailah trik pada pertanyaan 7.)
  \item Gambarlah \omterm{def:g1:addition:def}{jumlah} $1 + 2 + 3 + 4$ sebagai tangga persegi: satu
        persegi di kolom pertama, dua di kolom kedua, tiga di kolom
        ketiga, empat di kolom keempat. Tunjukkan pada gambarmu bahwa
        \emph{dua} tangga semacam itu, salah satunya dibalik ke
        bawah, cocok tepat menjadi satu persegi panjang dengan $4$
        baris berisi $5$ persegi --- lalu jelaskan mengapa ini adalah trik dua
        baris pada pertanyaan 7, yang digambar alih-alih ditulis.
  \item Sebuah jam berdentang $1$ kali pada pukul satu, $2$ kali pada
        pukul dua, \dots, $12$ kali pada pukul dua belas. Berapa
        dentang seluruhnya selama dua belas jam itu? Dan selama satu
        hari penuh?
  \item Satu lagi untuk Gauss: $5\,050$ detik --- apakah itu lebih
        atau kurang daripada satu setengah jam? Ubahlah menjadi jam,
        menit dan detik, dengan memakai pembagian Euklides
        (\cref{thm:g6:wholes:division}) dua kali.

\end{enumerate}
\end{problem}
